% Pursuing Stacks (cotes 134-2 to 134-8): archive cote 134-5, site page 88 (PDF page 89), typescript page 275.
% typescript. Draft compared: ps4.tex:4509-4587.
% AI-transcribed by Claude Opus 5.5 (claude-opus-5-5, adaptive thinking, effort high), Message Batches API, batch ps2 (msgbatch_01NVwr1Yuo8qJf5rykeGBXPK), 2026-11-12; reviewed against both supplied scan halves; see CHECK_DECISIONS.tsv.
\scanpage{134-5/88}
\mspage{\hand{275}}

\inserted[typed]{\hand{\uline{\textcolor{black}{Proposition.}}}} \xout{\hand{\uline{\textcolor{black}{Theorem}}}} Let $(M,M_a)$, $(M',M'_a)$ be two $\underline{W}$-asphericity structures ($\underline{W}$ a basic
localizer), $B \hand{\subset} M_a$ a full subcategory of $M$ generating the asphericity
structure $M_a$ , and
\[ f^* : M \hand{\longrightarrow} M' \]
a functor, admitting a left adjoint \hand{$f_!$} (hence $f^*$ commutes to inverse limits)
We consider the following six conditions, paralleling those of prop.4 of
section 75 (p.236):
\[ \text{(i)} \quad M_a = (f^*)^{-1}(M'_a ) , \]
\[ \text{(i')} \quad M_a \hand{\subset} (f^*)^{-1}(M'_a) , \quad \text{i.e.} \quad f^*(M_a) \hand{\subset} M'_a , \]
\[ \text{(i")} \quad B \hand{\subset} (f^*)^{-1}(M'_a) , \quad \text{i.e.} \quad f^*(B) \hand{\subset} M'_a , \]
\[ \text{(ii)} \quad W_a =(f^*)^{-1}(W'_a) , \]
\[ \text{(ii')} \quad W_a \hand{\subset} (f^*)^{-1}(W'_a) , \quad \text{i.e.} \quad f^*(W_a) \hand{\subset} W'_a , \]
\[ \text{(ii")} \quad \mathrm{F}\over{l}{\ell}(B) \hand{\subset} (f^*)^{-1}(W'_a) \quad \text{i.e.} \quad f^*(\mathrm{F}\over{l}{\ell}(B)) \hand{\subset} W'_a , \]
where $W_a$, $W'_a$ are the sets of $\underline{W}$-equivalences in $M$ and $M'$ respectively.
We'll make the extra assumption (I nearly forgot, sorry !) :
(Awk) There exists a \xout{test functor} $M'_a$-$\underline{W}$-aspheric functor
\[ i' : A \hand{\longrightarrow} M' \]
(A a small category), such that \inserted[typed]{$i =$} $f\hand{_!}i' : A \hand{\rightarrow} M$ factors through C.
Under these conditions and with these notations, the following holds:
The conditions (i),(i'),(i") are equivalent and imply the three others,
which satisfy the tautological implications (ii)$\hand{\Rightarrow}$(ii')$\hand{\Rightarrow}$(ii") . If
M and M' admit final objects and their asphericity structures are aspheric\hand{,}
then \xout{six} \inserted[typed]{the five} \inserted{first} conditions (i) to (ii') are equivalent, and if moreover $e_M \hand{\in} \mathrm{Ob}\, B$\hand{,}
all six are equivalent.

Proof: reduction to loc. cit. , using the commutative diagram
\margin{(1)}
% Hand-drawn arrows are marked blue; typed objects remain plain.
\diagram{\begin{tikzcd}[column sep=tiny]
M \ar[rr,"f^*",color=blue] \ar[dr,"i^*"',color=blue] \& \& M' \ar[dl,"i'^*",color=blue] \\
\& \hat{A} \&
\end{tikzcd}}
(up to isomorphism), and the relations
\[ M'_a =(i'^*)^{-1}(\hat{A}_a) \quad , \quad W'_a = (i'^*)^{-1}(\underline{W}_A) \quad . \]
This proof shows moreover that the conditions (i),(i'),(i") are equivalent
each to $i$ being $M_a$-$\underline{W}$-aspheric . Thus, if these conditions are satis-%
\margin{cf p.280\,$\times$}
fied (let's say then that $f^*$ is a \hand{\uline{\textcolor{black}{morphism}}}\inserted{(*)} for the asphericity structu\xout{xx}%
res ), we may use $i^{\hand{*}}$ and $i'^{\hand{*}}$ for describing the canonical functors from

\margin{88}
